% =============================================================================
% CompetitiveExam/CS301/10-file-organization-hashing-questions.tex
% Competitive Exam Practice 10 — GATE (Computer Science) Pattern
% File organization (sequential, direct, indexed sequential, hashed) and
% hashing: hash-function design, open hashing (separate chaining), closed
% hashing (linear/double probing), load factor, tombstones.
%
% Student-facing — NO answers. Answer key: 10-file-organization-hashing-answers.tex.
%
% Provenance (Tier 1: local GATEOverflow compilation CompetitiveExam/Books/
% GATE-PYQs/filter1_volume2.pdf, question text recovered by page-render + OCR;
% each answer independently re-derived by a Python trace this session and
% cross-checked against GeeksforGeeks / ExamSIDE / ExamGOAL mirrors):
%   Q1 [GATE CSE 1996 Q1.13, verified — GATEOverflow Vol.2 p.228]
%   Q2 [GATE CSE 2004 Q7, verified — GATEOverflow Vol.2 p.228]
%   Q3 [GATE CSE 2014 Set 1 Q40, verified — GATEOverflow Vol.2 p.230]
%   Q4 [GATE CSE 2007 Q40, verified — GATEOverflow Vol.2 p.229]
%   Q5 [GATE CSE 2020 Q23, verified — GATEOverflow Vol.2 p.39]
%   Q6 [GATE CSE 2023 Q10, verified — GATEOverflow Vol.2 p.54]
%   Q7-Q10 [Original, GATE-pattern] — fresh numbers, none reused from the deck
%     or the written assignment; every answer re-derived in the key.
%
% Compile with: cd CompetitiveExam/CS301 &&
%   TEXINPUTS=../../Preambles:$TEXINPUTS xelatex -interaction=nonstopmode 10-file-organization-hashing-questions.tex
% =============================================================================

\documentclass[11pt]{article}
\usepackage[margin=1in]{geometry}
% =============================================================================
% Preambles/header.tex — shared LaTeX/Beamer preamble for this project
%
% Use from a Beamer lecture by prepending:
%     \input{header}
% and compiling with TEXINPUTS=../Preambles:$TEXINPUTS (the /compile-latex
% skill does this automatically).
%
% PALETTE CONTRACT. Color names defined here MUST match the names in
% Quarto/theme-template.scss so Beamer and Quarto renderings use the same
% palette. The scripts/check-palette-sync.sh script enforces this.
%
% To customize: edit the HEX values below AND the matching $variable in
% Quarto/theme-template.scss, then run scripts/check-palette-sync.sh.
% =============================================================================

% -----------------------------------------------------------------------------
% Engine detection + encoding
%
% Our /compile-latex skill uses XeLaTeX, where inputenc is unnecessary and
% can produce encoding warnings. Only load inputenc under pdfTeX.
% -----------------------------------------------------------------------------
\usepackage{iftex}
\ifPDFTeX
  \usepackage[utf8]{inputenc}
\fi

\usepackage{amsmath, amssymb, amsthm}
\usepackage{booktabs}
\usepackage{graphicx}

% -----------------------------------------------------------------------------
% PALETTE — mirrors Quarto/theme-template.scss variable names.
% Load xcolor and define colors BEFORE anything that references them
% (notably hyperref, hypersetup, and the Beamer color assignments).
% -----------------------------------------------------------------------------
\usepackage{xcolor}

\definecolor{primary-blue}{HTML}{012169}      % headings, accents
\definecolor{primary-gold}{HTML}{B9975B}      % emphasis, borders
\definecolor{highlight-yellow}{HTML}{F2A900}  % markers, alerts
\definecolor{light-bg}{HTML}{E8EDF5}          % title / section backgrounds
\definecolor{jet}{HTML}{1A1A1A}               % body text

% Semantic colors (used in diagrams and annotations)
\definecolor{positive}{HTML}{15803D}          % good / observed / identified
\definecolor{negative}{HTML}{B91C1C}          % bad / problematic / bias
\definecolor{neutral}{HTML}{525252}           % reference / context

% Highlight accents (match .hi-* classes in the SCSS)
\definecolor{hi-slate}{HTML}{314F4F}
\definecolor{hi-green}{HTML}{2E8B57}
\definecolor{hi-red}{HTML}{C41E3A}

% Semantic aliases so skill templates and snippets can write
% \textcolor{accent}{...} without hard-coding "primary-blue".
\colorlet{accent}{primary-blue}
\colorlet{accent2}{primary-gold}

% -----------------------------------------------------------------------------
% Hyperref — loaded AFTER xcolor + palette so link colors resolve.
% -----------------------------------------------------------------------------
\usepackage{hyperref}
\hypersetup{colorlinks=true, linkcolor=primary-blue, urlcolor=primary-blue,
            citecolor=primary-blue, pdfborder={0 0 0}}

% -----------------------------------------------------------------------------
% Course code — set once per deck (after \input{header}, before \begin{document})
% via \coursecode{CS401}. Shown in the footer on every slide so a deck's course
% is identifiable without opening the title page. Empty by default.
% -----------------------------------------------------------------------------
\makeatletter
\newcommand{\coursecode}[1]{\gdef\@coursecodetext{#1}}
\gdef\@coursecodetext{}
\makeatother

% -----------------------------------------------------------------------------
% Beamer theme assignments (only apply under Beamer).
%
% We detect Beamer via \@ifclassloaded{beamer}, which is the documented,
% non-internal way. This must be wrapped in \makeatletter/\makeatother
% because \@ifclassloaded contains an @-catcode character.
% -----------------------------------------------------------------------------
\makeatletter
\@ifclassloaded{beamer}{%
  \usefonttheme{professionalfonts}
  \setbeamercolor{structure}{fg=primary-blue}
  \setbeamercolor{frametitle}{fg=primary-blue}
  \setbeamercolor{title}{fg=primary-blue}
  \setbeamercolor{subtitle}{fg=primary-gold}
  \setbeamercolor{author}{fg=primary-blue}
  \setbeamercolor{institute}{fg=neutral}
  \setbeamercolor{date}{fg=primary-gold}
  \setbeamercolor{itemize item}{fg=primary-blue}
  \setbeamercolor{itemize subitem}{fg=primary-gold}
  \setbeamercolor{alerted text}{fg=primary-gold}
  \setbeamercolor{block title}{fg=white, bg=primary-blue}
  \setbeamercolor{block body}{bg=light-bg}
  % Semantic block variants: block=definition/concept (blue), exampleblock=
  % worked example (green/positive), alertblock=key takeaway or caution (gold).
  % Keep to <=2 colored boxes per slide across ALL three variants (INV-7).
  \setbeamercolor{block title example}{fg=white, bg=positive}
  \setbeamercolor{block body example}{bg=light-bg}
  \setbeamercolor{block title alerted}{fg=white, bg=primary-gold}
  \setbeamercolor{block body alerted}{bg=light-bg}

  % Minimal footer: course code (left) + page X / N (right), no navigation clutter
  \setbeamertemplate{navigation symbols}{}
  \setbeamertemplate{footline}{%
    \color{neutral}\footnotesize\hspace{0.7em}\@coursecodetext\hfill
    \insertframenumber/\inserttotalframenumber\hspace{0.7em}\vspace{0.3em}}
}{}
\makeatother

% -----------------------------------------------------------------------------
% TikZ setup — libraries the snippet gallery uses
% -----------------------------------------------------------------------------
\usepackage{tikz}
\usetikzlibrary{arrows.meta, positioning, calc, decorations.pathreplacing,
                fit, shapes.geometric, backgrounds}

% Shared TikZ styles — snippets and hand-written diagrams can pick these up
% via `\begin{tikzpicture}[every node/.style={...}]` or by naming specific
% styles (e.g., `\node[dag-node] (X) at (0,0) {X};`).
\tikzset{
  % Node styles for causal DAGs and similar
  dag-node/.style = {
    draw=primary-blue, fill=light-bg, rounded corners=2pt,
    minimum width=1.2cm, minimum height=0.9cm, align=center, font=\small
  },
  decision-node/.style = {
    draw=primary-gold, fill=white, diamond, aspect=1.8,
    minimum width=1.8cm, minimum height=1.2cm, align=center, font=\small,
    inner sep=1pt
  },
  % Edge styles — observed vs counterfactual
  observed-edge/.style = {-{Latex[length=2.5mm]}, thick, draw=jet},
  counterfactual-edge/.style = {-{Latex[length=2.5mm]}, thick, draw=neutral, dashed},
  confound-edge/.style = {-{Latex[length=2.5mm]}, thick, draw=negative, dashed},
  % Data-point styles for plots
  observed-dot/.style = {fill=primary-blue, draw=primary-blue, circle, inner sep=1.2pt},
  counterfactual-dot/.style = {fill=white, draw=neutral, circle, inner sep=1.2pt}
}

% -----------------------------------------------------------------------------
% Convenience macros
% -----------------------------------------------------------------------------
\newcommand{\muted}[1]{\textcolor{neutral}{#1}}
\newcommand{\key}[1]{\textcolor{primary-gold}{\textbf{#1}}}
\newcommand{\good}[1]{\textcolor{positive}{#1}}
\newcommand{\bad}[1]{\textcolor{negative}{#1}}

% Standout frame macro: full-bleed dark background for section transitions.
% Beamer-only (uses \begin{frame}/\setbeamercolor/\usebeamerfont) -- do not
% call from an article-class file (e.g. Notes/<CODE>/*-notes.tex). \muted,
% \key, \good, \bad, and \coursecode above are plain xcolor/gdef and are
% safe to use from either class; verified when Notes/article-class support
% was added.
% Expand: \transitionslide{Your transition title}
\newcommand{\transitionslide}[1]{%
  \begingroup
  \setbeamercolor{background canvas}{bg=primary-blue}
  \begin{frame}[plain]
    \centering
    \vfill
    {\usebeamerfont{title}\color{white}\Large #1\par}
    \vfill
  \end{frame}
  \endgroup
}

% Section divider frame (Beamer-only), an alternative to \transitionslide with a
% darker two-line style: near-black (jet) background, a small white label line
% above a larger gold title, and a thin gold rule along the bottom edge.
% Expand: \sectiondivider{Section 2}{Pointers}
\newcommand{\sectiondivider}[2]{%
  \begingroup
  \setbeamercolor{background canvas}{bg=jet}
  \begin{frame}[plain]
    \vspace*{3.0cm}
    {\color{white}\ttfamily\large #1\par}
    \vspace{0.5cm}
    {\color{primary-gold}\LARGE #2\par}
    \begin{tikzpicture}[remember picture, overlay]
      \draw[primary-gold, line width=2.5pt]
        ($(current page.south west)+(0,0.1)$) -- ($(current page.south east)+(0,0.1)$);
    \end{tikzpicture}
  \end{frame}
  \endgroup
}

% -----------------------------------------------------------------------------
% Channel watermark — "Concepts That Click" (YouTube).
%
% Article-class files (CompetitiveExam Q/A sets, Notes, Assignments) get a
% light diagonal draft watermark on every page plus a centered footer naming
% the channel. Beamer decks get a subtle TikZ background watermark instead
% (the footline already carries the course code). Centralized here so every
% MCQ PDF is branded without per-file edits.
% -----------------------------------------------------------------------------
\makeatletter
\@ifclassloaded{article}{%
  \usepackage{draftwatermark}
  \SetWatermarkText{Concepts That Click}
  \SetWatermarkScale{1}
  \SetWatermarkFontSize{2cm}
  \SetWatermarkLightness{0.86}
  \SetWatermarkAngle{45}
  \usepackage{fancyhdr}
  \pagestyle{fancy}
  \fancyhf{}
  \fancyfoot[C]{\footnotesize\textcolor{neutral}{Concepts That Click~|~YouTube: \href{https://www.youtube.com/@concepts.thatclick}{@concepts.thatclick}}}
  \renewcommand{\headrulewidth}{0pt}
  \renewcommand{\footrulewidth}{0pt}
  \fancypagestyle{plain}{%
    \fancyhf{}%
    \fancyfoot[C]{\footnotesize\textcolor{neutral}{Concepts That Click~|~YouTube: \href{https://www.youtube.com/@concepts.thatclick}{@concepts.thatclick}}}%
    \renewcommand{\headrulewidth}{0pt}%
    \renewcommand{\footrulewidth}{0pt}%
  }
}{}
\@ifclassloaded{beamer}{%
  \setbeamertemplate{background}{%
    \begin{tikzpicture}[remember picture, overlay]
      \node[rotate=45, opacity=0.055, align=center]
        at (current page.center)
        {\Huge\textcolor{neutral}{Concepts That Click}};
    \end{tikzpicture}%
  }
}{}
\makeatother 
\coursecode{CS301}
\renewcommand{\thesection}{10.\arabic{section}}
\setlength{\emergencystretch}{2em}

\title{Competitive Exam Practice 10: File Organization and Hashing\\\large GATE (Computer Science) Pattern}
\author{Auqib Hamid Lone\\[0.3em]
  \emph{Department of Computer Science \& Engineering}, \emph{GCET Ganderbal}}
\date{\today}

\begin{document}
\maketitle

\noindent Each question carries 1--2 marks. MCQ = single correct option; NAT = Numerical Answer Type. Arrays and hash tables are 0-indexed unless a question states otherwise. ``Probe'' means one slot examined, counting the home slot as the first probe.

\begin{enumerate}
\item[\textbf{Q1}] \textit{[GATE CSE 1996 Q1.13, verified --- GATEOverflow Vol.~2 p.~228]} \textbf{[MCQ, 1 mark]} An advantage of a chained hash table (external hashing) over the open addressing scheme is
\begin{enumerate}
  \item[(A)] worst case complexity of search operations is less
  \item[(B)] space used is less
  \item[(C)] deletion is easier
  \item[(D)] none of the above
\end{enumerate}

\item[\textbf{Q2}] \textit{[GATE CSE 2004 Q7, verified --- GATEOverflow Vol.~2 p.~228]} \textbf{[MCQ, 2 marks]} Given the input
\[
(4322,\ 1334,\ 1471,\ 9679,\ 1989,\ 6171,\ 6173,\ 4199)
\]
and the hash function $h(x)=x \bmod 10$, which of the following statements are true?
\begin{enumerate}
  \item[I.] $9679$, $1989$, $4199$ hash to the same value.
  \item[II.] $1471$, $6171$ hash to the same value.
  \item[III.] All elements hash to the same value.
  \item[IV.] Each element hashes to a different value.
\end{enumerate}
\begin{enumerate}
  \item[(A)] I only
  \item[(B)] II only
  \item[(C)] I and II only
  \item[(D)] III or IV
\end{enumerate}

\item[\textbf{Q3}] \textit{[GATE CSE 2014 Set 1 Q40, verified --- GATEOverflow Vol.~2 p.~230]} \textbf{[MCQ, 2 marks]} Consider a hash table with $9$ slots. The hash function is $h(k)=k \bmod 9$. Collisions are resolved by chaining. The following $9$ keys are inserted in the order
\[
5,\ 28,\ 19,\ 15,\ 20,\ 33,\ 12,\ 17,\ 10.
\]
The maximum, minimum, and average chain lengths in the hash table, respectively, are
\begin{enumerate}
  \item[(A)] $3$, $0$, and $1$
  \item[(B)] $3$, $3$, and $3$
  \item[(C)] $4$, $0$, and $1$
  \item[(D)] $3$, $0$, and $2$
\end{enumerate}

\item[\textbf{Q4}] \textit{[GATE CSE 2007 Q40, verified --- GATEOverflow Vol.~2 p.~229]} \textbf{[MCQ, 2 marks]} Consider a hash table of size seven, with starting index zero, and a hash function $(3x+4) \bmod 7$. Assuming the table is initially empty, which of the following is the contents of the table when the sequence $1,3,8,10$ is inserted using closed hashing? A ``--'' denotes an empty location.
\begin{enumerate}
  \item[(A)] $8$, --, --, --, --, --, $10$
  \item[(B)] $1$, $8$, $10$, --, --, --, $3$
  \item[(C)] $1$, --, --, --, --, --, $3$
  \item[(D)] $1$, $10$, $8$, --, --, --, $3$
\end{enumerate}

\item[\textbf{Q5}] \textit{[GATE CSE 2020 Q23, verified --- GATEOverflow Vol.~2 p.~39]} \textbf{[NAT, 1 mark]} Consider a double hashing scheme in which the primary hash function is $h_1(k)=k \bmod 23$ and the secondary hash function is $h_2(k)=1+(k \bmod 19)$. The table size is $23$. The address returned by probe $1$ in the probe sequence (the sequence begins at probe $0$) for the key value $k=90$ is \underline{\hspace{1.5cm}}.

\item[\textbf{Q6}] \textit{[GATE CSE 2023 Q10, verified --- GATEOverflow Vol.~2 p.~54]} \textbf{[MCQ, 1 mark]} An algorithm has to store several keys generated by an adversary in a hash table. The adversary is malicious and tries to maximize the number of collisions. Let $k$ be the number of keys, $m$ the number of slots, and $k>m$. Which one of the following is the best hashing strategy to counteract the adversary?
\begin{enumerate}
  \item[(A)] Division method, i.e., use $h(k)=k \bmod m$.
  \item[(B)] Multiplication method, i.e., use $h(k)=\lfloor m(kA-\lfloor kA\rfloor)\rfloor$ for a carefully chosen constant $A$.
  \item[(C)] Universal hashing method.
  \item[(D)] If $k$ is a prime number, use the division method; otherwise use the multiplication method.
\end{enumerate}

\item[\textbf{Q7}] \textit{[Original, GATE-pattern]} \textbf{[Analysis, 2 marks]} A separate-chaining hash table has $m=2000$ slots and currently holds $n=1500$ keys.
\begin{enumerate}
  \item[(a)] Give the load factor $\alpha$.
  \item[(b)] Using the uniform-hashing approximations, give the expected number of key comparisons for a successful and for an unsuccessful search under chaining.
  \item[(c)] The same $n$ keys are stored in a table of $m=2000$ slots using linear probing, also at $\alpha=0.75$. Give the expected number of probes for a successful and for an unsuccessful search.
  \item[(d)] In one sentence, say which scheme is the safer choice at this load factor and why.
\end{enumerate}

\item[\textbf{Q8}] \textit{[Original, GATE-pattern]} \textbf{[Trace, 2 marks]} A table of $m=11$ slots uses $h(k)=k \bmod 11$ with linear probing. The keys
\[
12,\ 44,\ 13,\ 88,\ 23,\ 94
\]
are inserted in that order into an initially empty table.
\begin{enumerate}
  \item[(a)] Give the final table, with each key in its slot.
  \item[(b)] The key $13$ is deleted. Show the table after the deletion, using \texttt{DEL} for a tombstone.
  \item[(c)] After the deletion, how many probes does \texttt{search(88)} take, and why does it still succeed?
  \item[(d)] The key $35$ is then inserted. Which slot does it occupy, and how many probes does the insertion take?
  \item[(e)] Suppose slot $2$ had been cleared to \texttt{EMPTY} instead of marked \texttt{DEL}. What would \texttt{search(88)} return, and after how many probes?
\end{enumerate}

\item[\textbf{Q9}] \textit{[Original, GATE-pattern]} \textbf{[Design, 2 marks]} A university stores $400{,}000$ student records of $80$ bytes each. The disk block size is $4096$ bytes and records do not span blocks. The key is the roll number. The workload mixes point lookups (``find roll $X$'') with range queries (``list every roll in $[a,b]$'').
\begin{enumerate}
  \item[(a)] Give the number of records per block and the number of blocks in the file.
  \item[(b)] Choose one file organization and give the number of block reads for a point lookup, counting both the index and the data.
  \item[(c)] Give the number of block reads for a range query returning $5{,}000$ consecutive records under your choice.
  \item[(d)] What would the same range query cost in a hashed file, and why?
  \item[(e)] State one workload for which a direct (relative) file would beat both.
\end{enumerate}

\item[\textbf{Q10}] \textit{[Original, GATE-pattern]} \textbf{[Compute, 2 marks]} Compute each hash value, showing the intermediate value.
\begin{enumerate}
  \item[(a)] \emph{Mid-square}, $m=1000$: key $k=3214$. Give the middle four digits of $k^2$ and then the slot.
  \item[(b)] \emph{Folding}, $m=1000$: key $987654321$ split into three-digit groups from the right. Add the groups and reduce modulo $1000$.
  \item[(c)] \emph{Polynomial string hashing}, $m=97$, base $B=31$: give $h(\texttt{"dog"})$ using the character codes $d=100$, $o=111$, $g=103$.
  \item[(d)] In one sentence, why is a prime $m$ preferred for the division method $h(k)=k \bmod m$?
\end{enumerate}
\end{enumerate}

\end{document}
